\begin{table*}[tb]
\centering
\small
\begin{tabular}{l c}
\toprule
System & static-region area \\
\midrule
CausVid & 0.663 $\pm$ 0.073 \\
Reward-Forcing & 0.661 $\pm$ 0.060 \\
Self-Forcing & 0.651 $\pm$ 0.087 \\
Infinite-Forcing & 0.650 $\pm$ 0.090 \\
LongLive & 0.644 $\pm$ 0.095 \\
Rolling-Forcing & 0.616 $\pm$ 0.107 \\
Causal-Forcing & 0.614 $\pm$ 0.086 \\
\midrule
\multicolumn{2}{l}{\emph{Does mask area predict the drift score?}} \\
Spearman(area, fBD) over 7 systems & -0.75, exact $p=0.066$ \\
\midrule
\multicolumn{2}{l}{\emph{Ordering with precipitation excluded (3 of 23 prompts)}} \\
fBD ordering & best and worst unchanged \\
NBF ordering & identical \\
\bottomrule
\end{tabular}
\caption{\textbf{Partition and composition robustness.} The static region occupies a similar share of the frame for every audited system, so no system is scored over a much larger or smaller support than another. The area does correlate negatively with drift across systems; over seven systems that correlation is not determined, but its sign is the one circularity would produce---a system that drifts early enlarges its own dynamic region---and we flag it as an open validity question rather than dismissing it. Removing precipitation, the category a two-way partition cannot represent, leaves the normalised-background-flow ordering unchanged.}
\label{tab:robustness}
\end{table*}